\documentclass[10pt,a4paper]{amsart}
\usepackage[margin=19mm]{geometry}
\usepackage{amsmath,amssymb,mathtools}
\usepackage[scr=rsfs,cal=euler]{mathalfa}
\usepackage{xfrac}
\usepackage[unicode,hidelinks,bookmarks=false]{hyperref}
\newcommand{\ie}{i.e.\ }
\newcommand{\cf}{cf.\ }
\newcommand{\resp}{respectively\ }
\newcommand{\Subsection}{\subsection}
\providecommand{\nobreakdash}{\nobreak\hbox{-}}
\setlength{\emergencystretch}{2em}
\multlinegap0pt
\usepackage{amssymb}
\usepackage{xcolor}
\usepackage{algorithm}\usepackage{algpseudocode}
\usepackage{mathtools}
\usepackage{enumerate}
\usepackage{float}
\newtheorem{lemma}{Lemma}
\newcommand{\grad}{\mathrm{grad}}
\title{Riemannian conditional gradient methods for composite optimization problems}
\author{Kangming Chen, Ellen H. Fukuda}
\date{Source: arXiv:2412.19427v2 (CC BY 4.0)}
\begin{document}
\maketitle

\noindent\emph{Source and licence.} Riemannian conditional gradient methods for composite optimization problems. Version arXiv:2412.19427v2 (CC BY 4.0), distributed by arXiv under the Creative Commons Attribution 4.0 International licence, http://creativecommons.org/licenses/by/4.0/, as recorded in the arXiv licence metadata for this exact version and shown on the abstract page https://arxiv.org/abs/2412.19427v2. Full licence notice: \texttt{licenses/arxiv-2412.19427-CC-BY-4.0.txt}.\par\smallskip

\noindent\textit{Editorial scope.} Counted unit: the article's lemma lemma:armijo (source0.tex lines 615--620). The article's complete original proof of it is delivered with it (source0.tex lines 621--635, one \texttt{proof} environment of 15 lines). No mathematics is written, repaired or re-derived: the statement and the proof are the article's own lines, in the article's own notation, and no retained line is substituted or re-flowed.  Retained uncounted as the source-local context the proof needs: lines 412--414.  Nothing was omitted from any retained span, so the package carries no omission record.  No retained line contains a citation, so the wrapper carries no bibliography and no bibliography entry.  No retained line contains a figure environment, an image-inclusion command or drawing code, so no image is delivered and the package declares no asset dependency.  Editorial formatting only, in addition: an AMS \texttt{amsart} preamble that restates the article's own declarations and macros used by the retained lines, namely the environments \texttt{lemma} on separate counters, together with the standard packages those lines need.  The retained environments print in this document as Lemma 1 in order of appearance, whereas the article prints the same items with the numbering of its own full document.  The complete native source of the article as distributed in the arXiv e-print for this exact version is bundled unchanged as \texttt{sources/170758/source0.tex}.
    \begin{equation}\label{def theta}
    \theta(x^k) =  \left\langle\grad f(x^k), R^{-1}_{x^k}(p(x^k)) \right\rangle_{x^k} +g(p(x^k)) -g(x^k).
    \end{equation} 

\begin{lemma}\label{lemma:armijo}
    Let $\zeta \in(0,1)$, $x^k \in \mathcal{X}$ and $\theta(x^k) \neq 0$. Then, there exists $0<\bar{\lambda} \leq 1$ such that
    \begin{equation}\label{ineq armijo}
        F(R_{x^k} (\lambda R_{x^k}^{-1} (p(x^k)))) \leq F(x^k)-\zeta \lambda\left|\theta (x^k)\right|  \quad \forall \lambda \in(0, \bar{\lambda}] .
    \end{equation}
\end{lemma}

\begin{proof}
    Since $f$ is differentiable and $g$ is $R$-convex, the following holds for all $\lambda \in(0,1)$, 
    $$
    \begin{aligned}
    & F (R_{x^k} (\lambda R_{x^k}^{-1} (p(x^k)))) \\
    = & g (R_{x^k} (\lambda R_{x^k}^{-1} (p(x^k)))) + f(R_{x^k} (\lambda (R_{x^k}^{-1} (p(x^k))))) \\
    \leq & (1-\lambda) g (x^k) + \lambda g (p(x^k)) + f(x^k) + \lambda \left\langle \grad f (x^k), R_{x^k}^{-1} (p(x^k)) \right\rangle_{x^k} + {o}(\lambda) \\
    = & F (x^k) + \lambda \left( \left\langle \grad f (x^k), R_{x^k}^{-1} (p(x^k)) \right\rangle_{x^k} + g (p(x^k)) - g (x^k) \right) + {o}(\lambda) \\
    = & F(x^k) + \lambda \zeta \theta(x^k) + \lambda \left( (1-\zeta) \theta(x^k) + \frac{{o}(\lambda)}{\lambda} \right).
    \end{aligned}
    $$
    The inequality arises from the Taylor expansion of \(f\) and the $R$-convexity of \(g\).
    Since $\lim _{\lambda \rightarrow 0} \frac{{o}(\lambda)}{\lambda}=0$, $\zeta \in(0,1)$ and $\theta(x^k) < 0$ from the definition as in~\eqref{def theta}, there exists $\bar{\lambda} \in(0,1]$ such that $(1-\zeta) \theta(x^k)+\frac{{o}(\lambda)}{\lambda} \leq 0$,
    then \eqref{ineq armijo} holds for all $\lambda \in(0, \bar{\lambda}]$.
\end{proof}

\end{document}
